% 广义线性模型：讲解部分所用宏（按总笔记版式重新定义）
% 图形配色（沿用原名称，映射到总笔记色板）
\colorlet{InkSoft}{Muted}
\colorlet{PlotPrimary}{PKURed}
\colorlet{PlotSecondary}{black!65}
\colorlet{PlotThird}{black!40}
\colorlet{PlotFourth}{PKURed!45}
\colorlet{PlotNeutral}{PKUMuted}
\tikzset{every picture/.append style={line cap=round,line join=round},
  course observation/.style={PlotPrimary,line width=.75pt},
  course reference/.style={InkSoft,line width=.4pt,dashed}}
\pgfplotsset{every axis/.append style={axis line style={InkSoft,line width=.45pt},
  tick style={InkSoft,line width=.4pt},tick align=outside,label style={font=\small,color=Ink},
  tick label style={font=\footnotesize,color=InkSoft},title style={font=\small\bfseries\sffamily,color=Ink},
  legend style={draw=none,fill=white,font=\footnotesize,cells={anchor=west},inner sep=2pt},
  ymajorgrids=true,grid style={draw=black!9,line width=.2pt},axis background/.style={fill=white}}}
% 代码
\newcommand{\courseCJKMono}{}
\tcbuselibrary{listings}
\lstdefinestyle{course}{basicstyle=\ttfamily\fontsize{9.2}{13}\selectfont\color{Ink},
  columns=fullflexible,keepspaces=true,breaklines=true,breakatwhitespace=false,breakindent=14pt,
  showstringspaces=false,tabsize=2,numbers=left,numberstyle=\ttfamily\fontsize{7.3}{10}\selectfont\color{Muted!80},
  numbersep=9pt,stepnumber=1,firstnumber=1,xleftmargin=18pt,keywordstyle=\bfseries\color{CodeKeyword},
  keywordstyle=[2]\color{CodeFunction},keywordstyle=[3]\color{CodeConstant},stringstyle=\color{CodeString},
  commentstyle=\color{CodeComment}\itshape,aboveskip=0pt,belowskip=0pt}
\lstdefinelanguage{GLMR}{sensitive=true,morekeywords={if,else,for,while,repeat,function,return,next,break},morekeywords=[2]{glm,clogit,polr,multinom,coxph,cph,Surv,rcs,basehaz,cox.zph,summary,exp,coef,loglm,loglin,pchisq,c,list,library,require,install.packages,help,data,qcc.overdispersion.test,glm.nb,binomial,gaussian,poisson},morekeywords=[3]{TRUE,FALSE,NULL,NA,NaN,Inf},morecomment=[l]{\#},morestring=[b]',morestring=[b]",alsoletter={._},otherkeywords={<-}}
\lstdefinelanguage{GLMSAS}{sensitive=false,morekeywords={proc,run,model,class,by,direct,strata,weight,loglin,oddsratio},morekeywords=[2]{genmod,logistic,phreg,catmod,sort},morekeywords=[3]{data,dist,link,selection,ties,risklimits,cl,rl,lackfit,order,descending,ref,first,type3,obstats,residuals,type,scale,offset,SLE,SLS},morecomment=[l]{\#},morecomment=[s]{/*}{*/},morestring=[b]',morestring=[b]"}
\lstdefinelanguage{GLMStata}{sensitive=true,morekeywords={poisson,nbreg,stset,stcox,stcurve,estat,logistic,logit,clogit,mlogit,ologit},morekeywords=[2]{phtest},morekeywords=[3]{failure,survival,detail,r,irr,nolog,exposure,offset},morecomment=[l]{//},morecomment=[s]{/*}{*/},morecomment=[l]{\#},literate={\#\#}{{\textcolor{CodeKeyword}{\#\#}}}2,morestring=[b]',morestring=[b]"}
\lstdefinelanguage{GLMSPSS}{sensitive=false,morekeywords={COXREG,STATUS,STRATA,METHOD,PRINT,CRITERIA},morekeywords=[2]{ENTER,CI,PIN,POUT,ITERATE},morecomment=[l]{\#},morestring=[b]',morestring=[b]"}
\newtcblisting{coursecode}[1]{pkucode,listing only,listing engine=listings,
  listing options={style=course,language=GLM#1},title={#1},
  fonttitle=\sffamily\bfseries\small\color{black},coltitle=black,
  title after break={#1（续）}}
% 批注、推导、说明
\newenvironment{annotation}{\begin{zbnote}}{\end{zbnote}}
\newenvironment{derivation}[1]{\begin{zbderive}{#1}}{\end{zbderive}}
\newcommand{\revnote}[1]{\begin{zbcheck}#1\end{zbcheck}}
\newcommand{\sourcepages}[2]{}
% 表题、图题（编号随正文手写）
\newcommand{\ctable}[2]{\begin{center}\begin{minipage}{\linewidth}\centering\small
  \captionof*{table}{\sffamily\bfseries #1}#2\end{minipage}\end{center}}
\newcommand{\cfigure}[2]{\begin{center}\begin{minipage}{\linewidth}\centering #2
  \captionof*{figure}{\sffamily #1}\end{minipage}\end{center}}
% 记号
\DeclareMathOperator{\logit}{logit}
\let\Var\relax
\DeclareMathOperator{\Var}{Var}
\DeclareMathOperator{\SE}{SE}
\DeclareMathOperator{\OR}{OR}
\DeclareMathOperator{\RR}{RR}
\DeclareMathOperator{\HR}{HR}
\let\Cov\relax
\DeclareMathOperator{\Cov}{Cov}
\let\diag\relax
\DeclareMathOperator{\diag}{diag}
\DeclareMathOperator{\tr}{tr}
\DeclareMathOperator{\IRR}{IRR}
\renewcommand{\E}{\operatorname{E}}
\newcolumntype{L}[1]{>{\raggedright\arraybackslash}p{#1}}
